\begin{table}[htp]
	\centering
		\begin{tabular}{ | c | c | }
			\hline
			\rowcolor[HTML]{90D5FF} 
			\textbf{Hiperparámetro} & \textbf{Espacio de búsqueda}\\
			\hline
			
			activation (función de activación) & [relu, elu] \\
			\hline
			dropout (tasa de dropout)& [0.0, 0.1, 0.2, 0.3] \\
			\hline
			use\_batch\_norm (batch normalization) & [true, false] \\
			\hline
			optimizer & [sgd, adam] \\
			\hline
			lr (learning\_rate) & [0.0001, 0.01] \\
			\hline
			n\_hidden\_q (número capas ocultas para qgen7) & [1, 3] \\
			\hline
			n\_neurons\_q (neuronas por capa para qgen7) & [32, 64, 96, 128] \\
			\hline
			n\_hidden\_v (número capas ocultas para tensiones) & [1, 2] \\
			\hline
			n\_neurons\_v (neuronas por capa para tensiones) & [64, 96, 128, 160, 192, 224, 256] \\
			\hline
			n\_hidden\_p (número capas ocultas para pérdidas) & [2, 5] \\
			\hline
			n\_neurons\_p (neuronas por capa para pérdidas) & [64, 128, 192, 256] \\
			\hline
			
		\end{tabular} 
		
		\caption{Espacio de búsqueda de hiperparámetros (Red Neuronal). Fuente: Elaboración propia}
		\label{table:nn_hiperparam}
	\end{table}